\chapter{Geometric realisation and stability of quasi-geodesics}
\label{chap:stability}

Milestone M8. This chapter provides the tools for the equivalences of
Chapter~\ref{chap:hyperbolic} that pass between exact distances and quasi-isometries: the
geometric realisation of a graph, a geodesic space in which the vertices sit isometrically and
$\frac12$-densely; the four-point condition on a space with a quasi-dense isometric copy of a
space satisfying it; and the stability of quasi-geodesics in spaces with slim triangles
(Bridson--Haefliger, Theorem III.H.1.7). With them,
Theorem~\ref{Group.Generators.isQuasiHyperbolic_iff} replaces the word metric of a group by the
realisation of its Cayley graph (design decisions D18 and D28).

How the realisation is encoded in Lean is an open decision (D26); its nodes are not implemented
until it is decided. Stability is proved with chains of points instead of rectifiable paths (design
decision D27).

\section*{Geometric realisation of a graph}

\begin{definition}[Geometric realisation, cf.\ Bridson--Haefliger Def.~I.1.9]
  \label{SimpleGraph.Realization}
  Let $\Gamma$ be a connected simple graph with vertex set $V$ and graph distance $d_\Gamma$. Its
  \emph{geometric realisation} $|\Gamma|$ is obtained from a copy of $[0, 1]$ for each edge by
  gluing the endpoints of the copy for the edge $\{u, u'\}$ to the vertices $u$ and $u'$. A point
  $p$ of $|\Gamma|$ is a vertex or lies inside an edge; a point at parameter $s$ on an edge from $u$
  to $u'$ has distance $s$ from $u$ and $1 - s$ from $u'$. The distance between $p$ on an edge with
  endpoints $u, u'$ and $q$ on an edge with endpoints $v, v'$ is the least of the numbers
  \[ d(p, a) + d_\Gamma(a, b) + d(b, q), \qquad a \in \{u, u'\},\ b \in \{v, v'\}, \]
  and, if $p$ and $q$ lie on the same edge, of their difference of parameters. Each vertex $v$ is a
  point of $|\Gamma|$, which defines the \emph{vertex map} $V \to |\Gamma|$.
  \plannedlean{SimpleGraph.Realization.vertex}
\end{definition}

\paragraph{Differences from Bridson--Haefliger.} Bridson--Haefliger glue copies of $[0, 1]$ for
the edges of any graph with edge lengths $\lambda(e) > 0$, and define the distance as the infimum of
the lengths of piecewise linear paths; that this is a geodesic metric when there are finitely many
edge lengths is left to the reader (I.1.9, fact 2). Here every edge has length $1$, and the
distance is given by the formula above, which for unit edges is the length of a shortest path: a
shortest path from $p$ to $q$ either stays in one edge or leaves the edge of $p$ through an
endpoint $a$, follows a shortest walk from $a$ to the endpoint $b$ of the edge of $q$ through
which it arrives, and ends inside that edge. Only unit edges are needed (Cayley graphs), and the formula
reduces every fact to a finite case analysis on top of the graph distance. Whether the
construction follows the formula or the quotient is design decision D26.

\begin{lemma}[Vertices of the realisation]
  \label{SimpleGraph.Realization.dist_vertex}
  \uses{SimpleGraph.Realization, Metric.QuasiDenseRange}
  The vertex map of a connected graph $\Gamma$ preserves distances: the distance in $|\Gamma|$
  between two vertices is their graph distance. Every point of $|\Gamma|$ is within $\frac12$ of a
  vertex; in particular the vertex map has quasi-dense range.
  \plannedlean{SimpleGraph.Realization.quasiDenseRange_vertex}
\end{lemma}
\begin{proof}
  A vertex is an endpoint of each edge it lies on, at distance $0$ from itself, so the formula for
  two vertices $u, v$ has the term $d_\Gamma(u, v)$. Every other term is at least
  $d_\Gamma(u, v)$: a term through endpoints $a, b$ is at least $d_\Gamma(u, a) + d_\Gamma(a, b) +
  d_\Gamma(b, v)$, since a point at distance $1$ from a vertex along an edge is the other endpoint,
  adjacent to it. A point at parameter $s$ on an edge from $u$ to $u'$ is within $\min\{s, 1 - s\}
  \le \frac12$ of $u$ or $u'$.
\end{proof}

\begin{lemma}[The realisation is geodesic, cf.\ Bridson--Haefliger I.1.9]
  \label{SimpleGraph.Realization.isGeodesicSpace}
  \uses{SimpleGraph.Realization, Metric.IsGeodesic}
  The geometric realisation of a connected graph is a geodesic space.
\end{lemma}
\begin{proof}
  If $p, q$ lie on one edge and the minimum is their difference of parameters, the segment of the
  edge between them is a geodesic. Otherwise take endpoints $a, b$ attaining the minimum and a walk
  $a = w_0, \dots, w_n = b$ with $n = d_\Gamma(a, b)$. Run from $p$ to $a$ along its edge, then
  along the edges $w_i w_{i+1}$ at unit speed, then from $b$ to $q$. This map $\gamma$ is defined on
  $[0, d(p, q)]$, $1$-Lipschitz on each piece and hence everywhere, and joins $p$ to $q$. A
  $1$-Lipschitz map $\gamma \colon [0, \ell] \to X$ with $d(\gamma(0), \gamma(\ell)) = \ell$ is
  isometric: for $s \le t$, $d(\gamma(s), \gamma(t)) \ge \ell - d(\gamma(0), \gamma(s)) -
  d(\gamma(t), \gamma(\ell)) \ge \ell - s - (\ell - t) = t - s$.
\end{proof}

\section*{The four-point condition and quasi-dense subspaces}

The four-point condition passes to subspaces with the same constant
(\texttt{Metric.IsHyperbolicWith.of\_isometry} in the copied pull request). It also passes from a
quasi-dense subspace to the whole space, with a larger constant; this is how the word metric and the
realisation of the Cayley graph are compared.

\begin{lemma}[Four-point condition from a quasi-dense subspace]
  \label{Metric.IsHyperbolicWith.of_isometry_of_quasiDenseRange}
  \uses{Metric.IsHyperbolicWith, Metric.QuasiDenseRange}
  Let $f \colon X \to Y$ be an isometric map of pseudometric spaces such that every point of $Y$ is
  within $C$ of a point of $f(X)$. If $X$ satisfies the four-point condition with constant
  $\delta$, then $Y$ satisfies it with constant $\delta + 6C$.
\end{lemma}
\begin{proof}
  Given $w, x, y, z \in Y$, choose $w', x', y', z' \in X$ whose images are within $C$ of them.
  Each distance between two of the four points changes by at most $2C$ when they are replaced by the
  images, so each Gromov product, half a sum of three distances, changes by at most $3C$. The
  four-point condition for $w', x', y', z'$ in $X$, transported by the isometry, gives $(x \cdot
  z)_w \ge (fx' \cdot fz')_{fw'} - 3C \ge \min\{(fx' \cdot fy')_{fw'}, (fy' \cdot fz')_{fw'}\} -
  \delta - 3C \ge \min\{(x \cdot y)_w, (y \cdot z)_w\} - \delta - 6C$.
\end{proof}

\section*{Stability of quasi-geodesics}

Quasi-geodesics are defined in Chapter~\ref{chap:coarsely-geodesic}
(Definition~\ref{Metric.IsQuasiGeodesicWith}).

\inputleannode{Metric.HasSlimTrianglesWith.exists_hausdorffDist_le}
